% from ega4-2-fr.tex
\oldpage[IV-2]{20}
homéomorphisme sur son image, resp. un homéomorphisme bijectif), il en est de
même de $g$, et il s'agit donc d'une question purement topologique. Pour le
cas où $g\circ f$ est ouvert (resp. fermé), le fait que $g$ est alors ouvert
(resp. fermé) résulte de Bourbaki, \emph{Top. gén.}, chap.~I\textsuperscript{er},
3\textsuperscript{e} éd., \S~5, n\textsuperscript{o}~1, prop.~1~; pour les
deux autres cas, on peut se borner à supposer que
$g(f(X))=g(Y)=Z$, autrement dit au cas où $g\circ f$ est un homéomorphisme de
$X$ sur $Z$~; comme $f$ est surjectif, $g$ est nécessairement bijectif, et
comme $g$ est une application continue ouverte d'après ce qui précède, $g$
est bien un homéomorphisme de $Y$ sur $Z$.

Pour prouver (vi), notons que dire qu'un morphisme $g$ est ouvert (resp.
fermé, resp. un homéomorphisme sur son image, resp. un homéomorphisme
bijectif) équivaut à dire que $g_{\mathrm{red}}$ possède la même propriété.
D'autre part (I, 5.1.8), pour tout morphisme $Y'\to Y$, on a
\[
\bigl(X_{\mathrm{red}}\times_{Y_{\mathrm{red}}}Y'_{\mathrm{red}}\bigr)_{\mathrm{red}}
=(X\times_Y Y')_{\mathrm{red}},
\]
donc la remarque précédente montre que si $f_{\mathrm{red}}$ est
universellement ouvert (resp. universellement fermé, resp. universellement
bicontinu, resp. un homéomorphisme universel), il en est de même de $f$. La
réciproque se prouve de même, en notant ici que pour tout morphisme
$Y''\to Y_{\mathrm{red}}$, on a
\[
\bigl(X_{\mathrm{red}}\times_{Y_{\mathrm{red}}}Y''\bigr)_{\mathrm{red}}
=(X\times_Y Y'')_{\mathrm{red}}
\quad\text{(I, 5.1.8).}
\]

Enfin, la nécessité de (vii) résulte aussitôt de (iii). Inversement, supposons
remplie la condition (vii), et soit $g:Y'\to Y$ un morphisme~; alors les
$g^{-1}(U_\alpha)=U'_\alpha$ forment un recouvrement ouvert de $Y'$ et si
l'on note $f_\alpha$ la restriction $f^{-1}(U_\alpha)\to U_\alpha$ de $f$,
et $f'$ le morphisme $f_{(Y')}$, la restriction
$f'^{-1}(U'_\alpha)\to U'_\alpha$ n'est autre que
$(f_\alpha)_{(U'_\alpha)}$. On peut donc se borner à prouver que $f'$ est
ouvert (resp. fermé, resp. un homéomorphisme sur son image, resp. un
homéomorphisme sur $Y'$), ce qui est immédiat.

\begin{proposition}[2.4.4]
\label{IV.2.4.4-fr}
Un morphisme universellement bicontinu $f:X\to Y$ est radiciel (donc séparé
(I, 7.7.1)).
\end{proposition}

En effet, $f$ est universellement injectif par hypothèse (I, 3.5.11).

\begin{proposition}[2.4.5]
\EGAReaderTarget{ega4.c04.prop.2-4-5}{IV.2.4.5-fr}
\begin{enumerate}
\item[(i)] Un morphisme $f:X\to Y$ qui est entier, surjectif et radiciel est
un homéomorphisme universel.
\item[(ii)] Inversement, supposons $Y$ localement noethérien. Alors, si un
morphisme de type fini $f:X\to Y$ est un homéomorphisme universel, il est
fini, surjectif et radiciel.
\end{enumerate}
\end{proposition}

(i) Il suffit d'observer que les trois propriétés de $f$ se conservent par
changement de base (I, 3.5.2, I, 3.5.7 et II, 6.1.5), et comme un morphisme
entier est fermé (II, 6.1.10), il est clair que $f$ est un homéomorphisme de
$X$ sur $Y$.

(ii) Comme $f$ est de type fini et universellement fermé par hypothèse et
qu'il est séparé par (2.4.4), il est propre (II, 5.4.1), et pour tout
$y\in Y$, $f^{-1}(y)$ est réduit à un seul élément~; donc (III, 4.4.2) $f$
est fini~; il est clair que $f$ est surjectif, et il est radiciel puisqu'il
est universellement injectif (I, 3.5.11).

\begin{theorem}[2.4.6]
\label{IV.2.4.6-fr}
Soit $f:X\to Y$ un morphisme quasi-plat (2.3.3) et localement de présentation
finie. Alors $f$ est universellement ouvert. En particulier, un morphisme plat
localement de présentation finie est universellement ouvert.
\end{theorem}

On sait que pour tout changement de base $Y'\to Y$, $f_{(Y')}$ est
quasi-plat (2.3.3) et localement de présentation finie (1.4.3, (iii)). Il
suffit donc de prouver que $f$ est un
